\section*{Hoofdstuk \ref{ch:b3:facies-models} --- Faciesmodellen}

\begin{solution}{exo:b3:facies-models:1}
Van zee naar land: offshore, offshore-overgang met stormzand,
\omterm{def:b3:facies-models:shoreface}{onderwateroever}, strand, een kustvlakte met moerassen. Elk milieu kwam terecht
bovenop zijn zeewaartse buur: de kustlijn verplaatste zich zeewaarts, een regressie. Een
transgressie zou ze omgekeerd hebben opgestapeld, met de kustvlakte aan de
basis en offshore-slib bovenaan, vaak met een erosievlak waar de
\omterm{def:b3:facies-models:shoreface}{onderwateroever} voorbijtrok.
\end{solution}

\begin{solution}{exo:b3:facies-models:2}
(a) Een geulvulling of \omterm{def:b3:facies-models:rivers}{kronkelwaard}: een grindige uitschuring aan de basis, kruisgelaagdheid,
daarna rimpels en wortelhoudend slib. (b) Een prograderende \omterm{def:b3:facies-models:shoreface}{onderwateroever} of deltafront:
slib met stormzandlagen en \omterm{def:b3:facies-models:hcs}{hummocky kruisgelaagdheid} eronder,
kruisgelaagd zand van de \omterm{def:b3:facies-models:shoreface}{onderwateroever} en strand erboven. (c) Opgestapeld zand van \omterm{def:b3:facies-models:rivers}{verwilderde rivieren},
een eolische duin of carbonaatbank: kruisgelaagde sets met constante korrelgrootte,
met grote duinfronten bij het eolische geval, of skeletale en oolitische
korrels voor de bank.
\end{solution}

\begin{solution}{exo:b3:facies-models:3}
$\phi = (2.65 - 2.40)/(2.65 - 1.0) = 0.152$; met pekel,
$0.25/1.55 = 0.161$.
\end{solution}

\begin{solution}{exo:b3:facies-models:4}
$e^{-2\pi z/\lambda} = e^{-\pi} = 0.043$ bij $z = \lambda/2$: ongeveer
\qty{38}{m} en \qty{55}{m}. Bij \qty{20}{m}: $e^{-2\pi\times20/75} = 0.19$
en $e^{-2\pi\times20/110} = 0.32$ van de oppervlaktebeweging.
\end{solution}

\begin{solution}{exo:b3:facies-models:5}
$(60 - 47.6)/(189 - 47.6) = 0.088$; $(60 - 52.6)/(189 - 52.6) = 0.054$. Een
matrixwaarde met een onzekerheid van \qty{5}{\micro\second/ft} verandert
de porositeit met een derde: de time-average-vergelijking heeft een matrixwaarde nodig die op kernen is gecontroleerd.
\end{solution}

\begin{solution}{exo:b3:facies-models:6}
$(6.0 + 3.5 + 8.0 + 2.5)/40 = 20/40 = 0.50$. Poriënvolume: $20\times0.22
\times10^6 = \qty{4.4e6}{m^3}$ per vierkante kilometer.
\end{solution}

\begin{solution}{exo:b3:facies-models:7}
Dicht bij de canyon: alleen Ta, de bovenste divisies verwijderd door de volgende stroom;
langs de waaier: Tc (rimpels) en Td (silt) onder Te (slib). De stroming verliest snelheid
en haar grove vracht met de afstand; alleen de fijne fractie reist
\qty{30}{km}, en wordt afgezet door een langzamere stroming die rimpels opbouwt, niet
door neerslag uit een dichte suspensie.
\end{solution}

\begin{solution}{exo:b3:facies-models:8}
Een \omterm{def:b3:facies-models:rivers}{kronkelwaard} van een \omterm{def:b3:facies-models:rivers}{meanderende rivier}: lag op de uitgeschuurde geulbodem,
kruisgelaagd zand van de diepe geul, zand met rimpels bovenop de bar,
slib van de overstromingsvlakte met wortels. De opeenvolging legt de geul vast van bodem tot
oever, dus de geul was ongeveer \qty{9}{m} diep (iets minder, na
compactie van het slib). Niet over de basis: het is een erosievlak, en
wat eronder lag kan elke oudere afzetting zijn geweest.
\end{solution}

\begin{solution}{exo:b3:facies-models:9}
$60/40\,000 = 0.0015$, \ang{0.09}: een ramp. $300/1500 = 0.20$, \ang{11.3},
na een vlakke top van \qty{5}{m} diep: een \omterm{def:b3:facies-models:ramp}{omzoomd plat}, met rif- of bankkorrels
aan de rand en blokken en korrelstromen op de helling eronder.
\end{solution}

\begin{solution}{exo:b3:facies-models:10}
Het gordel heeft tijd $w/v$ nodig om een vaste plaats te passeren, waarin de bodem
stijgt ten opzichte van het zeeniveau met $A$ eronder: het zand vult $h = A\,w/v$.
$2000\times0.2/20 = \qty{20}{m}$. Bij $A = 0$ laat de \omterm{def:b3:facies-models:shoreface}{onderwateroever} alleen
de dikte van het eigen profiel achter, een dunne laag, omdat de gordels over elkaar schuiven
zonder aggradatie. De verhouding $A/v$ is de tangens van de hoek van het
\omterm{def:b3:sequence-stratigraphy:trajectory}{kustlijntraject} van de propositie in Hoofdstuk~\ref{ch:b3:sequence-stratigraphy}: de
dikte van elke faciës is de breedte maal die tangens.
\end{solution}

\begin{solution}{exo:b3:facies-models:11}
Dichtheid: $(2.65 - 2.30)/1.65 = 0.212$; sonisch: $(86 - 52.6)/136.4 = 0.245$.
Met $\phi = 0.245$, $\rho_{\mathrm{ma}} = (2.30 - 0.245)/(1 - 0.245) =
\qty{2.72}{g/cm^3}$: dichter dan kwarts, zoals voor een zandsteen gecementeerd door
calciet (2.71) of die calcietkorrels bevat.
\end{solution}

\begin{solution}{exo:b3:facies-models:12}
Offshore-slib op een kolenlaag brengt niet-aangrenzende milieus in contact: een
\omterm{def:b3:sequence-stratigraphy:parasequence}{overstromingsvlak}, hier een \omterm{def:b3:sequence-stratigraphy:surfaces}{transgressief vlak} door golferosie, het schelpgruis
een lag. Toen de kustlijn landwaarts verschoof, sneden de golven de bodem
tot ongeveer de \omterm{def:b3:facies-models:shoreface}{golfbasis bij mooi weer}: het strand, de bovenste \omterm{def:b3:facies-models:shoreface}{onderwateroever} en
de top van de kustvlakte, tot ongeveer \qty{10}{m} van de sectie, kunnen
ontbreken.
\end{solution}

\begin{solution}{pb:b3:facies-models:1}
\textbf{1.} Krijt (koud, lage dichtheid), schalie vanaf \qty{766}{m} (heet), kalksteen vanaf
\qty{866}{m} (koud, dicht, snel), zandsteen vanaf \qty{1046}{m} (klokvormige
pieken), schist vanaf \qty{1166}{m}.

\textbf{2.} Krijt \qty{766}{m}, schalie 100, kalksteen 180, zandsteen
\qty{120}{m}; basis van de deklaag op $441.0 - 1166 = \qty{-725}{m}$.

\textbf{3.} $(2.71 - 2.03)/1.71 = 0.40$ nabij het oppervlak, $(2.71 -
2.36)/1.71 = 0.20$ op \qty{750}{m}: het krijt verliest de helft van zijn porositeit
onder het gewicht van de bovenliggende lagen, door compactie.

\textbf{4.} De korrels zijn kleien van onzekere dichtheid, en het water is
deels aan hen gebonden: de aannamen van de vergelijking gaan niet op.

\textbf{5.} Kristallijn basement: de deklaag wordt doorboord, en geen reservoir wordt
onder die laag verwacht.

\textbf{6.} $(2.71 - 2.573)/1.71 = 0.080$.

\textbf{7.} $(58.9 - 47.6)/141.4 = 0.080$: hetzelfde.

\textbf{8.} $(2.65 - 2.32)/1.65 = 0.200$.

\textbf{9.} $(79.9 - 52.6)/136.4 = 0.200$: hetzelfde.

\textbf{10.} $0.33/1.55 = 0.213$, één procentpunt en een derde meer.

\textbf{11.} Het zandsteen, met een porositeit van \qty{20}{\percent} in zijn zand;
de schalie boven de kalksteen dicht het af, en de overstromingsmodders verdelen het
zandsteen in compartimenten.

\textbf{12.} Een klokvorm: lage gammawaarde in het kruisgelaagde zand, die oploopt door het fijner wordende interval
naar de hoge waarden van de wortelmodder, op een scherpe basis.

\textbf{13.} Een \omterm{def:b3:facies-models:association}{opwaarts fijner wordende opeenvolging}: geulbodemresidu op een erosieve
basis, geuldunen (kruisgelaagdheid), top van de \omterm{def:b3:facies-models:rivers}{kronkelwaard} (rimpels), overstromingsvlakte
(wortelmodder).

\textbf{14.} De geul, zijn \omterm{def:b3:facies-models:rivers}{kronkelwaard} op de binnenoever, oeverwallen en een
overstromingsvlakte met bodems, naast elkaar.

\textbf{15.} Ongeveer de dikte van zand en fijner wordend interval, $10.4 + 3.2 =
\qty{13.6}{m}$, iets meer vóór compactie.

\textbf{16.} Een \omterm{def:b3:facies-models:rivers}{meanderende rivier} waarvan de geulgordel herhaaldelijk verschoof
door \omterm{def:b3:facies-models:rivers}{avulsie} terwijl het bekken zonk, waarbij elke nieuwe geul insnijdt in de
overstromingsvlakte die de vorige achterliet.

\textbf{17.} Nee: elke basis is erosief; de overstromingsmodder eronder en het
geulzand erboven hoeven niet naast elkaar gelegen te hebben.

\textbf{18.} $(60 - 30)/(105 - 30) = 0.40$.

\textbf{19.} Laagste: $14\times(0.60 + 0.40\times0.25) = \qty{9.80}{m}$;
dikste: $16\times(0.65 + 0.40\times0.20) = \qty{11.68}{m}$.

\textbf{20.} Schoon zand \qty{58.8}{m}; fijner wordende intervallen \qty{33.4}{m}; modder
\qty{27.8}{m}.

\textbf{21.} Netto $= 58.8 + 0.40\times33.4 = \qty{72.2}{m}$; 
\omterm{def:b3:facies-models:ntg}{net-to-gross-verhouding} $72.2/120 = 0.60$.

\textbf{22.} Fracties 0.20 en 0.60: \qty{65.5}{m} en \qty{78.8}{m}, een
spreiding van \qty{13}{m} door alleen al de keuze van de afkapwaarde.

\textbf{23.} $72.2\times0.20\times10^6 = \qty{1.44e7}{m^3}$ per
vierkante kilometer.

\textbf{24.} De netto zanddikte van het Trias-zandsteen bij Varden-2 is
\textbf{\qty[text-series-to-math=true]{72.2}{m}}.
\end{solution}
